\chapter{Измерение, унитарная эволюция и квантовая динамика}

\lectureoverview{Сначала мы соберём уже используемые принципы измерения: наблюдаемые — эрмитовы операторы, их собственные значения являются возможными результатами, различимые состояния ортогональны, а правило Борна задаёт вероятности. Новая проблема — динамика. Обратимость требует единой линейной эволюции, сохраняющей перекрытия; отсюда следует, что оператор эволюции унитарен, а его бесконечно малый генератор эрмитов. Этим генератором является гамильтониан, и его действие даёт зависящее от времени уравнение Schr\"odinger. Затем мы выводим формулу для среднего значения, коммутаторное уравнение для средних, его связь с классическими скобками Пуассона и закон сохранения среднего значения энергии.}

Пока оставим спиновую систему в стороне. Мы вернёмся к ней после того, как будет выстроена общая структура. Принципы ниже не предлагаются как минимальный независимый набор. Некоторые из них подкрепляют другие или частично влекут их, но отдельная формулировка каждого облегчает понимание их физического содержания.

\section{Какие результаты может дать измерение}

\emph{Наблюдаемая} — это то, что можно измерить. Устоявшийся термин несколько возвышеннее, чем необходимо; \emph{измеримое} подошло бы ничуть не хуже. Каждой наблюдаемой мы сопоставляем эрмитов оператор. Пусть пока \(L\) обозначает такой оператор:
\begin{equation}
L^\dagger=L.
\end{equation}
Не всякий линейный оператор эрмитов. И наоборот, вопрос о том, можно ли измерить данный эрмитов оператор с помощью имеющихся материалов и аппаратуры, может оказаться трудной экспериментальной задачей. Тем не менее абстрактное рабочее правило достаточно щедро: эрмитов оператор представляет возможную наблюдаемую, даже если практический способ её измерения ещё не изобретён.

Если
\begin{equation}
L\ket{\lambda}=\lambda\ket{\lambda},
\end{equation}
то \(\lambda\) — возможный численный результат измерения \(L\). Мы используем один и тот же символ для обозначения собственного значения и соответствующего ему собственного вектора. Запись в виде кета различает их: \(\lambda\) — число, а \(\ket{\lambda}\) — состояние.

Почему именно эрмитовы операторы? Простейший прибор сообщает вещественные числа. Если прибор как будто сообщает одно комплексное число, то в действительности он сообщает две независимые вещественные величины. Поэтому собственные значения наблюдаемой должны быть вещественными. Более того, состояния, соответствующие различным определённым результатам, должны быть физически различимы и, следовательно, ортогональны. В используемой здесь конечномерной ситуации линейный оператор с полным ортонормированным базисом собственных векторов и вещественными собственными значениями является эрмитовым. Таким образом, физические утверждения и эрмитова структура согласуются друг с другом.

\subsection{Различимость — это ортогональность}

Выражение \emph{физически различимы} имеет операциональный смысл. Предположим, что спин подготовлен либо вверх, либо вниз. Измерение вдоль оси \(z\) разделяет эти возможности без неоднозначности: вверх даёт \(+1\), вниз даёт \(-1\). Точно так же состояния влево и вправо можно различить, измеряя вдоль оси \(x\).

Теперь предположим, что подготовлено либо состояние вверх, либо состояние вправо. Измерение \(z\) даёт \(+1\) с достоверностью для состояния вверх, но состояние вправо тоже даёт \(+1\) с ненулевой вероятностью. Измерение \(x\) лишь обращает эту асимметрию. Наклон прибора на \(45^\circ\) не устраняет перекрытия между распределениями результатов. Ни одна ориентация не даёт результата, который всегда позволял бы различить эти два варианта подготовки.

\begin{definition}
Два состояния физически различимы, если некоторый эксперимент или набор экспериментов может однозначно определить, какой вариант подготовки был реализован.
\end{definition}

Квантовая механика представляет полную различимость геометрически:
\begin{equation}
\text{distinguishable states}
\quad\Longrightarrow\quad
\text{orthogonal vectors}.
\end{equation}

\begin{classroomqa}
\audiencequestion{Может ли другая ось измерения различить состояние вверх и состояние вправо?}
\lecturerresponse{Нет. Измерение вдоль \(z\) или \(x\) иногда может установить, какая была подготовка, но по крайней мере один результат остаётся совместимым с обеими. Для промежуточной оси вероятности результатов тоже перекрываются. Препятствие связано не с неудачным выбором прибора; два вектора состояния не ортогональны.}
\end{classroomqa}

\section{Правило Борна}

Собственные значения говорят нам, что может произойти, но не как часто. Именно правило вероятностей придаёт предсказательное содержание. Если нормированная система подготовлена в состоянии \(\ket{A}\), а \(\ket{\lambda}\) — нормированный собственный вектор измеряемой наблюдаемой \(L\), то
\begin{equation}
P(\lambda)=\left|\braket{\lambda}{A}\right|^2.
\label{eq:l4-born-rule}
\end{equation}
Перекрытие \(\braket{\lambda}{A}\) есть \emph{амплитуда вероятности}. Вообще говоря, оно комплексно, так что само по себе вероятностью быть не может. Его модуль в квадрате вещественен и неотрицателен.

\begin{figure}[htbp]
\centering
\includegraphics[width=0.86\textwidth]{lecture_04_figure_01.png}
\caption{Правило Борна, записанное рядом с принципами измерения: вероятность есть модуль перекрытия в квадрате.}
\lectureframe{00:12:40}
\end{figure}

Выберем ортонормированный базис из собственных векторов \(L\) и разложим
\begin{equation}
\ket{A}=\sum_i \alpha_i\ket{\lambda_i},
\qquad
\alpha_i=\braket{\lambda_i}{A}.
\end{equation}
Тогда правило Борна принимает вид
\begin{equation}
P(\lambda_i)=|\alpha_i|^2.
\end{equation}
Нормировка затем даёт
\begin{equation}
\sum_i P(\lambda_i)=\sum_i|\alpha_i|^2=1.
\end{equation}
Требования положительности и нормировки сами по себе не выводят правило Борна, но они показывают, почему модуль в квадрате так естественно согласуется с теорией вероятностей. В конечном счёте это правило проверяется экспериментом.

\begin{classroomqa}
\audiencequestion{Требуется ли модуль в квадрате также для того, чтобы все вероятности складывались в единицу?}
\lecturerresponse{Для нормированного состояния, разложенного по ортонормированному базису, квадраты модулей амплитуд действительно дают в сумме единицу. Поэтому положительность и нормировка полностью согласуются с правилом Борна. Какое из этих утверждений следует считать логически первичным, здесь менее важно, чем их взаимная согласованность; правило вероятностей устанавливается экспериментом.}
\end{classroomqa}

\begin{classroomqa}
\audiencequestion{Почему собственное значение, которое является числом, записывают внутри кета?}
\lecturerresponse{Обозначение \(\ket{\lambda}\) означает собственный вектор, соответствующий численному собственному значению \(\lambda\). В более педантичной записи можно было бы ввести другую метку, но слишком большое число индексов быстро затемняет вычисление. Символ кета сохраняет различие между числом и вектором.}
\end{classroomqa}

\begin{classroomqa}
\audiencequestion{Какой принцип превращает несколько возможных результатов в один фактический результат?}
\lecturerresponse{Вероятности задаёт правило Борна. Если подготовленное состояние уже является собственным вектором \(L\), то соответствующий результат имеет вероятность единица. Для суперпозиции реализуется один из результатов с вероятностями из уравнения ~\eqref{eq:l4-born-rule}. Связанное с этим правило обновления состояния, или правило коллапса, — отдельный вопрос измерения; к нему мы вернёмся после введения запутанности.}
\end{classroomqa}

\subsection{Вырожденные результаты}

Одному собственному значению может соответствовать более одного независимого собственного вектора. Предположим,
\begin{align}
L\ket{\lambda,1}&=\lambda\ket{\lambda,1},\\
L\ket{\lambda,2}&=\lambda\ket{\lambda,2}.
\end{align}
Любая линейная комбинация в их линейной оболочке также является собственным вектором с собственным значением \(\lambda\):
\begin{equation}
L\bigl(\alpha\ket{\lambda,1}+\beta\ket{\lambda,2}\bigr)
=
\lambda\bigl(\alpha\ket{\lambda,1}+\beta\ket{\lambda,2}\bigr).
\end{equation}
Эта свобода позволяет выбрать ортонормированный базис внутри вырожденного собственного подпространства. Если метка \(a\) пробегает этот базис, то вероятность численного результата \(\lambda\) равна сумме по его взаимно исключающим ортогональным каналам:
\begin{equation}
P(\lambda)=\sum_a\left|\braket{\lambda,a}{A}\right|^2.
\label{eq:l4-degenerate-born}
\end{equation}
Эквивалентно, если \(P_\lambda\) — проектор на всё собственное подпространство, то \(P(\lambda)=\bra{A}P_\lambda\ket{A}\). Обозначение через проектор компактно; уравнение ~\eqref{eq:l4-degenerate-born} непосредственно показывает сложение вероятностей.

\begin{classroomqa}
\audiencequestion{Что, если подготовленное состояние перекрывается с несколькими собственными векторами или нескольким собственным векторам соответствует одно и то же собственное значение?}
\lecturerresponse{Для разных собственных значений каждое перекрытие даёт вероятность своего результата. Для нескольких ортогональных собственных векторов с одним и тем же собственным значением нужно возвести модули перекрытий в квадрат и сложить их. Если векторы, первоначально выбранные в вырожденном подпространстве, не ортогональны, сначала замените их ортонормированным базисом; суммарная вероятность от выбора базиса не зависит.}
\end{classroomqa}

\begin{classroomqa}
\audiencequestion{Можно ли получить однозначный ответ для спина, подготовка которого полностью неизвестна?}
\lecturerresponse{Не по одному произвольному экземпляру. Без ограничений на способ его подготовки не существует одного измерения, которое однозначно определяло бы неизвестное квантовое состояние. Информация о подготовке задаёт альтернативы, которые измерение может различить.}
\end{classroomqa}

\section{Новый вопрос: как меняются состояния}

Предыдущие принципы задают состояния и измерения. Динамика спрашивает, что происходит между одним моментом времени и другим. Классическая аналогия полезна. Монета с двумя состояниями может оставаться неизменной на каждом такте, а может чередовать орёл, решка, орёл, решка. Для шестигранного кубика возможны более разнообразные правила обновления, а для непрерывного пространства состояний — бесконечно многие. Динамический закон отображает текущее состояние в следующее.

Удовлетворительный классический закон детерминистичен и обратим. Различные состояния не сливаются в одно, потому что иначе настоящее перестало бы определять прошлое. Такое сохранение информации иногда называют \emph{минус первым законом}: он предшествует более привычным законам сохранения.

Квантовая версия начинается с вектора состояния \(\ket{\psi(t)}\), который явно рассматривается как функция времени. Постулируем единый линейный оператор временной эволюции \(U(t)\) такой, что
\begin{equation}
\ket{\psi(t)}=U(t)\ket{\psi(0)}.
\label{eq:l4-time-development}
\end{equation}
Один и тот же \(U(t)\) действует на любое начальное состояние системы. Он не выбирается заново в зависимости от входного вектора, и его действие линейно. Можно вообразить и более общие нелинейные или зависящие от состояния законы, но они быстро вступают в противоречие с вероятностной и причинной структурой квантовой механики, а также со структурой составных систем.

\subsection{Обратимость сохраняет перекрытия}

Пусть второе начальное состояние эволюционирует по тому же закону:
\begin{equation}
\ket{\phi(t)}=U(t)\ket{\phi(0)}.
\end{equation}
Если эти два приготовления изначально различимы, они должны оставаться различимыми. Поэтому
\begin{equation}
\braket{\psi(0)}{\phi(0)}=0
\quad\Longrightarrow\quad
\braket{\psi(t)}{\phi(t)}=0.
\end{equation}

\begin{figure}[htbp]
\centering
\includegraphics[width=0.86\textwidth]{lecture_04_figure_02.png}
\caption{Два состояния эволюционируют под действием одного и того же \(U(t)\); начальная ортогональность сохраняется.}
\lectureframe{00:34:43}
\end{figure}

Сохранения нормированных состояний и ортогональности достаточно, чтобы сохранялись все скалярные произведения. Чтобы увидеть механизм, разложим два произвольных вектора по ортонормированному базису. Базисные векторы остаются ортонормированными при эволюции, поэтому коэффициенты восстанавливают то же самое перекрытие в любой момент времени:
\begin{equation}
\braket{\psi(t)}{\phi(t)}=\braket{\psi(0)}{\phi(0)}.
\label{eq:l4-overlap-conservation}
\end{equation}
Тем самым сохраняется не только полная различимость, но и вся степень сходства между двумя состояниями.

Бра-вектор, соответствующий уравнению~\eqref{eq:l4-time-development}, получается эрмитовым сопряжением:
\begin{equation}
\bra{\psi(t)}=\bra{\psi(0)}U^\dagger(t).
\end{equation}
Подстановка эволюции бра- и кет-векторов в уравнение~\eqref{eq:l4-overlap-conservation} даёт
\begin{equation}
\bra{\psi(0)}U^\dagger(t)U(t)\ket{\phi(0)}
=
\braket{\psi(0)}{\phi(0)}
\end{equation}
для любой пары начальных векторов. Если оператор имеет между любой парой векторов те же матричные элементы, что и единичный оператор, то он и есть единичный оператор. Следовательно
\begin{equation}
U^\dagger(t)U(t)=I.
\label{eq:l4-unitarity}
\end{equation}
Оператор, удовлетворяющий уравнению~\eqref{eq:l4-unitarity}, является \emph{унитарным}. В конечномерном случае
\begin{equation}
U^\dagger(t)=U^{-1}(t),
\end{equation}
так что эволюцию можно обратить.

\begin{classroomqa}
\audiencequestion{Обозначают ли здесь прописная и строчная формы \(\Phi\) разные виды состояний?}
\lecturerresponse{Нет. Это лишь непоследовательность в записи. Мы последовательно используем \(\ket{\phi}\); математика различает состояния по их меткам и значениям, а не по случайному использованию прописных букв.}
\end{classroomqa}

\begin{classroomqa}
\audiencequestion{Можно ли читать \(U^\dagger U=I\) как утверждение о том, что один множитель является обратным к другому?}
\lecturerresponse{Да. Для унитарного оператора обратный оператор равен \(U^{-1}=U^\dagger\). Это алгебраическое выражение обратимой эволюции, сохраняющей перекрытия.}
\end{classroomqa}

\section{Бесконечно малый интервал времени и гамильтониан}

При нулевом прошедшем времени ничего не меняется, поэтому \(U(0)=I\). Для очень малого интервала \(\epsilon\) оператор \(U(\epsilon)\) должен быть близок к единичному. В единицах, где временно полагают \(\hbar=1\), запишем
\begin{equation}
U(\epsilon)=I-i\epsilon H+O(\epsilon^2).
\label{eq:l4-small-u}
\end{equation}
Знак минус и множитель \(i\) на первых порах являются соглашениями; не зная \(H\), их можно поглотить его определением. Их полезность проявляется, когда накладывается условие унитарности. Эрмитово сопряжение даёт
\begin{equation}
U^\dagger(\epsilon)=I+i\epsilon H^\dagger+O(\epsilon^2).
\end{equation}
Тогда
\begin{align}
I
&=U^\dagger(\epsilon)U(\epsilon)\\
&=\bigl(I+i\epsilon H^\dagger\bigr)
  \bigl(I-i\epsilon H\bigr)+O(\epsilon^2)\\
&=I+i\epsilon(H^\dagger-H)+O(\epsilon^2).
\end{align}
Член первого порядка должен обращаться в нуль:
\begin{equation}
H^\dagger=H.
\label{eq:l4-h-hermitian}
\end{equation}

\begin{figure}[htbp]
\centering
\includegraphics[width=0.86\textwidth]{lecture_04_figure_03.png}
\caption{Расчёт унитарности в первом порядке требует, чтобы бесконечно малый генератор был эрмитовым.}
\lectureframe{00:43:38}
\end{figure}

Буква \(H\) означает не Hermitian, а \emph{гамильтониан}. Поскольку он эрмитов, он представляет наблюдаемую. Его отождествляют с энергией и, на более структурном уровне, с генератором временных трансляций. Это отождествление не исчерпывается названием: в системах с известным классическим пределом квантовый \(H\) должен воспроизводить динамику классического гамильтониана. Симметрия, степени свободы, физическое суждение и, в конечном счёте, эксперимент определяют, какой эрмитов оператор является гамильтонианом конкретной системы.

\begin{classroomqa}
\audiencequestion{Можно ли использовать в качестве \(H\) любой эрмитов оператор? Что определяет, какой из них является фактическим?}
\lecturerresponse{Пока что вывод допускает любой эрмитов оператор в гильбертовом пространстве системы. Физика выбирает один, построенный из степеней свободы этой системы. Для спина он должен быть построен из спиновых операторов и внешних полей. Симметрия и известные взаимодействия сужают возможности; эксперимент фиксирует оставшиеся константы, а иногда и сам вид оператора.}
\end{classroomqa}

Гамильтониан может явно зависеть от времени. Простой пример даёт магнитное поле, создаваемое током, который изменяют. Если внешние параметры не зависят от времени, то \(H\) не имеет явной зависимости от времени; это инвариантность относительно временных трансляций. Если параметры принудительно изменяются во времени, энергия не обязана сохраняться.

\begin{classroomqa}
\audiencequestion{В разложении по малому времени является ли \(H\) константой, или должно быть \(H(0)\)?}
\lecturerresponse{Для не зависящей от времени задачи \(H\) — один и тот же оператор в любой момент времени. В более общем случае локальный шаг, начинающийся в момент \(t\), содержит \(H(t)\). Здесь вывод сначала рассматривает случай инвариантности относительно временных трансляций, а затем отмечает, как явная зависимость от времени меняет закон сохранения энергии.}
\end{classroomqa}

\section{Зависящее от времени уравнение Schr\"odinger}

Применим уравнение~\eqref{eq:l4-small-u} к состоянию:
\begin{equation}
\ket{\psi(\epsilon)}=(I-i\epsilon H)\ket{\psi(0)}+O(\epsilon^2).
\end{equation}
Переставляя члены и деля на \(\epsilon\),
\begin{equation}
\frac{\ket{\psi(\epsilon)}-\ket{\psi(0)}}{\epsilon}
=-iH\ket{\psi(0)}+O(\epsilon).
\label{eq:l4-difference-quotient}
\end{equation}

\begin{figure}[htbp]
\centering
\includegraphics[width=0.86\textwidth]{lecture_04_figure_04.png}
\caption{Бесконечно малое обновление состояния, переписанное в виде векторнозначного разностного отношения.}
\lectureframe{00:51:42}
\end{figure}

В дифференцировании вектора нет никакой особой загадки. Два вектора можно вычесть, их разность можно разделить на число, и предел определяет ещё один вектор. Выбранный выше в качестве нуля момент отсчёта может быть любым мгновением. Поэтому, переходя к пределу \(\epsilon\to0\), получаем
\begin{equation}
\boxed{\frac{d}{dt}\ket{\psi(t)}=-iH\ket{\psi(t)}}
\qquad(\hbar=1).
\label{eq:l4-schrodinger-natural}
\end{equation}

\begin{figure}[htbp]
\centering
\includegraphics[width=0.86\textwidth]{lecture_04_figure_05.png}
\caption{Завершённый вывод на доске приводит к заключённому в рамку зависящему от времени уравнению Schr\"odinger.}
\lectureframe{00:56:11}
\end{figure}

Это обобщённое зависящее от времени уравнение Schr\"odinger. Знакомое дифференциальное уравнение для волновой функции частицы — одно из представлений этого абстрактного уравнения для вектора состояния. Чтобы им пользоваться, нужно знать \(H\) и начальное состояние.

Для конечного интервала нужно решить дифференциальное уравнение. Когда \(H\) не зависит от времени, многократное составление бесконечно малых шагов даёт
\begin{equation}
U(t)=e^{-iHt}
\qquad(\hbar=1),
\end{equation}
это и есть точная версия составления всех малых шагов. Для гамильтонианов, зависящих от времени, требуется временное упорядочение; это уточнение мы пока откладываем.

\begin{classroomqa}
\audiencequestion{В разложении использовался бесконечно малый интервал. Что определяет состояние спустя долгое время?}
\lecturerresponse{Многократно применяйте бесконечно малое обновление или, что то же самое, решайте дифференциальное уравнение. Для не зависящего от времени \(H\) произведение малых шагов превращается в \(e^{-iHt/\hbar}\).}
\end{classroomqa}

\subsection{Восстановление постоянной Планка}

В предыдущем уравнении использовались единицы, в которых \(\hbar=1\). Размерностный анализ показывает, где должно стоять \(\hbar\). Производная имеет размерность, обратную времени, тогда как \(H\) имеет размерность энергии. Поскольку \(\hbar\) имеет размерность произведения энергии на время,
\begin{equation}
\boxed{i\hbar\frac{d}{dt}\ket{\psi(t)}=H\ket{\psi(t)}}.
\label{eq:l4-schrodinger}
\end{equation}
Эквивалентно, бесконечно малый оператор равен \(U(\epsilon)=I-i\epsilon H/\hbar+O(\epsilon^2)\), а конечная не зависящая от времени эволюция имеет вид \(U(t)=e^{-iHt/\hbar}\).

\begin{classroomqa}
\audiencequestion{Где именно входит постоянная Планка?}
\lecturerresponse{Постоянная Планка входит в любой из двух эквивалентных форм
\begin{align*}
i\hbar\,\partial_t\ket{\psi}&=H\ket{\psi},\\
\partial_t\ket{\psi}&=-\frac{i}{\hbar}H\ket{\psi}.
\end{align*}
Обе записи имеют согласованные размерности, поскольку \(\hbar\) имеет размерность произведения энергии на время.}
\end{classroomqa}

Состояние эволюционирует детерминистически согласно уравнению~\eqref{eq:l4-schrodinger}; результаты измерений остаются статистическими, потому что по правилу Борна они выбираются из непрерывно изменяющегося состояния.

\begin{classroomqa}
\audiencequestion{Зависит ли вывод от обратимости, и как он может сохраняться при измерении, которое возмущает состояние?}
\lecturerresponse{Он зависит от линейной эволюции и сохранения скалярных произведений. Уравнение управляет изолированной системой. Измерительный прибор взаимодействует достаточно сильно, чтобы стать частью физической системы, и им нельзя пренебречь. Объединённую систему и прибор нужно рассматривать квантовомеханически; их запутанность — именно то, что необходимо понять, прежде чем можно будет анализировать измерение.}
\end{classroomqa}

\section{Математическое ожидание — это среднее}

Чтобы сравнить квантовую динамику с классическими уравнениями, нам нужно среднее значение наблюдаемой. Стандартный термин \emph{математическое ожидание} потенциально вводит в заблуждение. Для случайной величины, с равной вероятностью принимающей только значения \(+1\) и \(-1\), среднее равно нулю, хотя нуль никогда не наблюдается. Симметричное распределение может даже обращаться в нуль во всей окрестности своего среднего. Среднее, медиана и мода совпадают для некоторых гладких одновершинных распределений, но не в общем случае.

Если \(L\) может принимать значения \(\lambda_i\) с вероятностями \(P(\lambda_i)\), определим
\begin{equation}
\overline{\lambda}=\sum_i\lambda_iP(\lambda_i).
\label{eq:l4-probability-average}
\end{equation}
Это первый момент распределения, а не обещание относительно результата одного испытания.

\begin{classroomqa}
\audiencequestion{Разве не следует ожидать именно наиболее часто встречающееся значение?}
\lecturerresponse{Наиболее часто встречающееся значение — это мода; взвешенное среднее — это среднее значение. Для простого узкого распределения они могут совпадать, но у двугорбого распределения среднее может лежать в области практически нулевой вероятности. Слово «ожидание» относится к математическому среднему, а не обязательно к наиболее вероятному отдельному исходу.}
\end{classroomqa}

Повторные испытания связывают математическое среднее с данными. При обычных предположениях о независимости и конечности моментов законы больших чисел утверждают сходимость по вероятности выборочного среднего. Они не гарантируют, что всякая конечная последовательность окажется близкой к среднему; чрезвычайно длинная серия одинаковых результатов измерения спина остаётся возможной, лишь маловероятной. Это та же статистическая тонкость, которая не позволяет второму закону термодинамики быть микроскопическим запретом без исключений.

\begin{classroomqa}
\audiencequestion{Если в честном эксперименте со спином случайно выпадает миллион «вверх» подряд, то его среднее и в самом деле равно «вверх»?}
\lecturerresponse{Выборочное среднее для этой конкретной серии действительно равно «вверх». Теоретическое математическое ожидание остаётся нулевым. По мере роста числа независимых испытаний выборочное среднее с высокой вероятностью приближается к теоретическому среднему, но теория вероятностей никогда не исключает все исключительные конечные последовательности.}
\end{classroomqa}

\subsection{Операторный сэндвич}

Разложим нормированное состояние по ортонормированному собственному базису \(L\):
\begin{equation}
\ket{A}=\sum_i\alpha_i\ket{\lambda_i},
\qquad
\bra{A}=\sum_j\alpha_j^*\bra{\lambda_j}.
\end{equation}
Из правила Борна и уравнения~\eqref{eq:l4-probability-average} получаем
\begin{equation}
\overline{\lambda}=\sum_i\lambda_i|\alpha_i|^2.
\end{equation}
Теперь вычислим
\begin{align}
\bra{A}L\ket{A}
&=\sum_{i,j}\alpha_j^*\alpha_i
  \bra{\lambda_j}L\ket{\lambda_i}\\
&=\sum_{i,j}\alpha_j^*\alpha_i\lambda_i
  \braket{\lambda_j}{\lambda_i}\\
&=\sum_i\lambda_i|\alpha_i|^2.
\end{align}
Двойная сумма схлопывается, поскольку \(\braket{\lambda_j}{\lambda_i}=\delta_{ji}\). Следовательно,
\begin{equation}
\boxed{\expect{L}_A=\bra{A}L\ket{A}}.
\label{eq:l4-expectation-sandwich}
\end{equation}

\begin{figure}[htbp]
\centering
\includegraphics[width=0.86\textwidth]{lecture_04_figure_06.png}
\caption{Разложение по собственному базису сводит вероятностное среднее к операторному сэндвичу \(\bra{A}L\ket{A}\).}
\lectureframe{01:11:50}
\end{figure}

Для средних будут использоваться и угловые скобки, и черта сверху. Когда состояние эволюционирует,
\begin{equation}
\expect{L}_t=\bra{\psi(t)}L\ket{\psi(t)}.
\end{equation}
Нижний индекс лишь напоминает, что состояние, а значит и среднее, зависит от времени. В настоящем обсуждении в картине Schr\"odinger сам \(L\) фиксирован, если явно не указана зависимость от времени.

\begin{classroomqa}
\audiencequestion{Следует ли писать \(t\) вне скобок, и означает ли черта сверху выборочное среднее?}
\lecturerresponse{Любое обозначение можно снабдить пометой, показывающей зависимость от времени, но стандартное квантовое обозначение — \(\expect{L}_t\), или просто \(\expect{L}\), где время понятно из \(\ket{\psi(t)}\). Здесь черта сверху обозначает теоретическое вероятностное среднее, первый момент \(\sum_i\lambda_iP(\lambda_i)\), а не специально выборочное среднее по конечному набору данных.}
\end{classroomqa}

\section{Как эволюционируют средние}

Продифференцируем этот сэндвич по правилу произведения:
\begin{equation}
\frac{d}{dt}\expect{L}
=
\left(\frac{d}{dt}\bra{\psi}\right)L\ket{\psi}
+
\bra{\psi}L\left(\frac{d}{dt}\ket{\psi}\right).
\label{eq:l4-product-derivative}
\end{equation}
В единицах \(\hbar=1\) уравнение для кета и его эрмитово-сопряжённое имеют вид
\begin{align}
\frac{d}{dt}\ket{\psi}&=-iH\ket{\psi},\\
\frac{d}{dt}\bra{\psi}&=+i\bra{\psi}H.
\end{align}
При комплексном сопряжении знак \(i\) меняется, и \(H^\dagger=H\). Подстановка в уравнение~\eqref{eq:l4-product-derivative} даёт
\begin{align}
\frac{d}{dt}\expect{L}
&=i\bra{\psi}HL\ket{\psi}
-i\bra{\psi}LH\ket{\psi}\\
&=i\bra{\psi}(HL-LH)\ket{\psi}.
\end{align}

Операторы не обязаны коммутировать. Их порядок существен, так же как существен порядок матриц. Определим
\begin{equation}
[H,L]=HL-LH.
\end{equation}
Тогда
\begin{equation}
\boxed{\frac{d}{dt}\expect{L}=i\expect{[H,L]}}
\qquad(\hbar=1).
\label{eq:l4-average-evolution-natural}
\end{equation}

\begin{figure}[htbp]
\centering
\includegraphics[width=0.86\textwidth]{lecture_04_figure_07.png}
\caption{Производная среднего по времени определяется средним значением коммутатора.}
\lectureframe{01:22:22}
\end{figure}

Восстанавливая размерности,
\begin{equation}
\boxed{\frac{d}{dt}\expect{L}
=\frac{i}{\hbar}\expect{[H,L]}
=-\frac{i}{\hbar}\expect{[L,H]}}.
\label{eq:l4-average-evolution}
\end{equation}
Если сам \(L\) явно зависит от времени, то общее уравнение содержит дополнительный член \(\expect{\partial L/\partial t}\). Здесь \(L\) фиксирован.

\begin{classroomqa}
\audiencequestion{В уравнении с коммутатором действительно стоит плюс \(i\)?}
\lecturerresponse{Да; при \([H,L]=HL-LH\) получается \(+(i/\hbar)\expect{[H,L]}\). Если переставить аргументы коммутатора, его знак меняется на противоположный, что даёт эквивалентную форму \(-(i/\hbar)\expect{[L,H]}\).}
\end{classroomqa}

\begin{classroomqa}
\audiencequestion{Уравнение для бра было получено комплексным сопряжением или эрмитовым сопряжением, и почему производная остаётся слева?}
\lecturerresponse{Это эрмитово сопряжение. Производная кета есть кет, а двойственный вектор этой производной есть производная соответствующего бра. Это можно увидеть непосредственно из разностного отношения: разности бра — это бра, точно так же как разности кетов — это кеты. Время здесь — параметр, а не оператор, порядок которого нужно обращать.}
\end{classroomqa}

\section{Классическая тень}

В классической гамильтоновой механике функция фазового пространства \(L(q,p)\), не имеющая явной зависимости от времени, подчиняется соотношению
\begin{equation}
\frac{dL}{dt}=\{L,H\}_{\mathrm{PB}},
\label{eq:l4-poisson-evolution}
\end{equation}
где \(\{\cdot,\cdot\}_{\mathrm{PB}}\) — скобка Пуассона. Сравните это с уравнением~\eqref{eq:l4-average-evolution}. Обе скобки антисимметричны:
\begin{align}
\{L,H\}_{\mathrm{PB}}&=-\{H,L\}_{\mathrm{PB}},\\
[L,H]&=-[H,L].
\end{align}
Обе также действуют на произведение как производные:
\begin{align}
\{LM,H\}_{\mathrm{PB}}
&=\{L,H\}_{\mathrm{PB}}M+L\{M,H\}_{\mathrm{PB}},\\
[LM,H]
&=L[M,H]+[L,H]M.
\label{eq:l4-commutator-product}
\end{align}
Тождество для коммутатора получается, если прибавить и вычесть \(LHM\):
\begin{align}
LMH-HLM
&=(LMH-LHM)\\
&\quad +(LHM-HLM)\\
&=L(MH-HM)\\
&\quad +(LH-HL)M.
\end{align}

Эта параллель подсказывает соответствие
\begin{equation}
[A,B]
\quad\longleftrightarrow\quad
i\hbar\{A,B\}_{\mathrm{PB}},
\label{eq:l4-dirac-correspondence}
\end{equation}
или, эквивалентно,
\begin{equation}
\{A,B\}_{\mathrm{PB}}
\quad\longleftrightarrow\quad
-\frac{i}{\hbar}[A,B].
\end{equation}

\begin{figure}[htbp]
\centering
\includegraphics[width=0.86\textwidth]{lecture_04_figure_08.png}
\caption{Классическое уравнение со скобкой Пуассона и квантовое соответствие для коммутатора с восстановленным \(\hbar\).}
\lectureframe{01:32:32}
\end{figure}

Дирак распознал эту связь и использовал её как путь в квантовую механику. Тем не менее обращаться с ней нужно как с соответствием, а не как с универсальным правилом подстановки: могут возникать неоднозначности порядка операторов, а классический предел должен выводиться из квантовой динамики. Концептуально квантовая механика первична; классическая механика является её приближением для подходящих состояний и масштабов. Коммутаторы первичны по отношению к скобкам Пуассона, хотя историческая аналогия помогает обнаружить эту связь.

\begin{classroomqa}
\audiencequestion{Почему \(H\) называют энергией и почему делают вывод о сохранении энергии ещё до введения принципа симметрии?}
\lecturerresponse{На данном этапе \(H\) строго является эрмитовым генератором сдвигов по времени. Называть его энергией — это физическое отождествление, которое нужно проверить там, где известен классический предел. Для больших или тяжёлых систем в подходящих состояниях квантовые средние должны следовать классическим уравнениям, а квантовый \(H\) должен соответствовать классическому гамильтониану. Приведённое ниже утверждение о сохранении предполагает, что этот генератор не имеет явной зависимости от времени.}
\end{classroomqa}

\section{Средняя энергия сохраняется}

Положим \(L=H\) в уравнении~\eqref{eq:l4-average-evolution} при не зависящем от времени гамильтониане:
\begin{equation}
\frac{d}{dt}\expect{H}
=\frac{i}{\hbar}\expect{[H,H]}.
\end{equation}
Всякий оператор коммутирует сам с собой,
\begin{equation}
[H,H]=HH-HH=0,
\end{equation}
поэтому
\begin{equation}
\boxed{\frac{d}{dt}\expect{H}=0}.
\label{eq:l4-energy-conservation}
\end{equation}

\begin{figure}[htbp]
\centering
\includegraphics[width=0.86\textwidth]{lecture_04_figure_09.png}
\caption{Сохранение средней энергии следует непосредственно из того, что гамильтониан коммутирует сам с собой.}
\lectureframe{01:39:30}
\end{figure}

Это доказывает намеренно скромное утверждение: средняя энергия постоянна. Более сильное утверждение обо всём распределении энергии будет следовать из той же динамики. Если \(H\) явно зависит от времени, дополнительная явная производная не обязана обращаться в нуль, и вынужденная система может обмениваться энергией с тем, что изменяет её параметры.

Остаются открытыми две темы. Вычисление вероятности для спина, измеряемого вдоль произвольной оси, всё ещё ждёт своего разбора, и задача обновления состояния ещё не выведена. Между вмешательствами изолированное состояние эволюционирует унитарно. Во время измерения система и прибор взаимодействуют и запутываются.

\begin{classroomqa}
\audiencequestion{Куда делось обещанное вычисление вероятности для спина, измеряемого вдоль произвольной оси?}
\lecturerresponse{Общая динамика заняла всё доступное время. Расчёт для спина остаётся следующим конкретным применением, и к нему вернёмся после изложения этой структурной схемы.}
\end{classroomqa}

\begin{classroomqa}
\audiencequestion{Как постулаты могут быть полными без постулата коллапса?}
\lecturerresponse{Пока что «коллапс» — лишь предварительный способ говорить. Измерение устанавливает запутанность между системой и прибором. Пока составные системы и запутанность не будут разработаны, процесс измерения нельзя моделировать честно. Это и есть следующий концептуальный шаг.}
\end{classroomqa}